\scanpage{144}
% archive sheet number: 143
\mspage{2.}

\uline{Theorem 2.3}: \uline{If} $A$ \uline{is the valuation ring in} $k$ \uline{and if} $M$ \uline{is a} $G_A$-\uline{module such that} $mA=A$, \uline{then} $\operatorname{Card}H^0(A,M)=\operatorname{Card}H^1(A,M)$ \uline{and} $H^2(A,M)=0$. \uline{In the duality between} $H^1(k,M)$ \uline{and} $H^1(k,\unsure{\overset{\infty}{M}|\widetilde M})$ \uline{the subgroups} $H^1(A,M)$ \uline{and} $H^1(A,\overset{\infty}{M})$ \uline{are the annihilator of each other}.

\uline{Remark}: Theorem 2.1 has been generalized by S. Shatz (Harvard Thesis, June 1962) to the case where $M$ is an arbitrary Artinian \inserted{commutative} group scheme defined over $k$, and where we take cohomology with respect to arbitrary finite extensions, not only the separable ones. \deleted{\unsure{We}\ \illegible{2}} The cohomology groups are no longer finite, but can be endowed with a canonical locally compact topology. For example if $\operatorname{char}k=p>0$, then using the exact sequences

\diagram{\begin{tikzcd}[ampersand replacement=\&,column sep=small]
0\arrow[r]\&\mu_p\arrow[r]\&\unsure{\mathbb G|G}_m\arrow[r,"p"]\&\mathbb G_m\arrow[r]\&0
\end{tikzcd},\quad
\begin{tikzcd}[ampersand replacement=\&,column sep=small]
0\arrow[r]\&\alpha_p\arrow[r]\&\mathbb G_a\arrow[r,"\unsure{P|p}"]\&\mathbb G_a\arrow[r]\&0
\end{tikzcd},\quad
\begin{tikzcd}[ampersand replacement=\&,column sep=small]
0\arrow[r]\&\mathbb Z/p\mathbb Z\arrow[r]\&\mathbb G_a\arrow[r,"\wp"]\&\mathbb G_a\arrow[r]\&0
\end{tikzcd}}

we find
\diagram{\[\begin{array}{lll}
H^0(k,\mu_p)=0&H^0(k,\alpha_p)=0&H^0(k,\mathbb Z/p\mathbb Z)=\mathbb Z/p\mathbb Z\\
H^1(k,\mu_p)=k^*/k^{*p}\ \text{(compact)}&
H^1(k,\alpha_p)=k/k^p\ \text{(locally compact)}&
H^1(k,\mathbb Z/p\mathbb Z)=k/\wp k\ \text{(discrete)}\\
H^2(k,\mu_p)=\frac1p\mathbb Z/\mathbb Z&H^2(k,\alpha_p)=0&H^2(k,\mathbb Z/p\mathbb Z)=0
\end{array}\]}

and in these cases the dualities corresponding to $\mu_p=\widetilde{\mathbb Z/p\mathbb Z}$ and $\alpha_p=\widetilde\alpha_p$ are well known, being made explicit by $\langle x,y\rangle=\operatorname{Tr}\operatorname{res}(x\frac{dy}{y})$ and $=\operatorname{Tr}\operatorname{res}(x\,dy)$ respectively.

\S\,3. \uline{Global results}: Let now $k$ be a finite extension of $\mathbb Q$ (case $(N)$) or a function field in one variable over a finite field (case $(F)$). Let $S$ be a non-empty (possibly infinite) set of prime divisors of $k$, including the archimedean primes in case $(N)$; and let $A$ be the ring of elements in $k$ which are integers at all primes \uline{not} in $S$. (For example, if $S$ is the set of \uline{all} prime divisors, then $A=k$.)

\uline{Theorem 3.1}: \uline{There is a natural exact sequence of continuous maps}

\margin{$(\overset{\infty}{M})$}
\diagram{\begin{tikzcd}[ampersand replacement=\&,column sep=small,row sep=large,scale=.58,transform shape]
0\arrow[r]\&H^0(A,\overset{\infty}{M})\arrow[r,"\alpha_0"]\&
\prod_{\unsure{P|p}\in S}H^0(k_p,\overset{\infty}{M})\arrow[r,"\beta_0"]\&
\widehat{H^2(A,M)}\arrow[r,"\gamma"]\&H^1(A,\overset{\infty}{M})\arrow[d,bend left=30,"\alpha_1"]\\
\&\&\begin{gathered}\text{\uline{and isomorphisms}}\\
H^r(A,M)\xrightarrow[\sim]{\alpha_r}\prod_{P\in S}H^{\unsure{0|r}}(k_p,M)\ \text{\uline{for}}\ r\geq3.
\end{gathered}\&\&
\mathop{\unsure{\underline{\prod}|\prod}}\limits_{P\in S}H^1(k_p,\overset{\infty}{M})\arrow[d,bend left=30,"\beta_1"]\\
0\&\widehat{H^0(A,M)}\arrow[l]\&
\coprod_{P\in S}H^2(k_p,\overset{\infty}{M})\arrow[l,"\beta_2"]\&
H^2(A,\overset{\infty}{M})\arrow[l,"\alpha_2"]\&
\widehat{H^1(A,M)}\arrow[l,"\delta"]
\end{tikzcd}}

\uline{where the symbols are explained as follows}: $k_p$ denotes the completion
